Aunque el conjunto de datos original ya incluía un total de 83 características por flujo, estas describen propiedades estadísticas, temporales y del tamaño del flujo. Por este motivo, se introdujeron características adicionales que describen el contenido de la carga útil del flujo, es decir, características asociadas al \textit{payload}.

La primera de estas es la entropía de Shannon. Esta mide el grado de aleatoriedad o impredecibilidad de una secuencia de símbolos. Un valor bajo indica un contenido más regular y predecible, mientras que un valor alto indica un contenido más disperso y menos predecible. Sea $X$ una variable aleatoria discreta, la expresión general de la entropía de Shannon es:

\begin{equation}

H(X) = - \sum_{i=1}^{n} p(x_i)\log_2\bigl(p(x_i)\bigr)
\end{equation}

donde $p(x_i)$ es la probabilidad de aparición del símbolo $x_i$.

Si se aplica esta definición al problema, se puede deducir que $X \in \{0,1,2,\dots,255\}$, ya que un byte puede tomar 256 valores. En este caso, la probabilidad de aparición de un byte no es una probabilidad exacta, debido a que no se puede saber con certeza la probabilidad de que aparezca ese valor en tráfico real. Por este motivo, la probabilidad de aparición de un byte será una probabilidad estimada, que se obtiene de la siguiente forma:

\begin{equation}
\label{EQ:PROBABILIDAD_BYTE}
p_i = \frac{n_i}{N}
\end{equation}

donde $n_i$ es el número de veces que aparece el byte $i$ en el flujo y $N$ es el número total de bytes observados.

Por lo tanto, la expresión de la entropía de Shannon adaptada a este problema queda:

\begin{equation}
\label{EQ:SHANNON_BYTES}
H(X) = - \sum_{i=0}^{255} \frac{n_i}{N}\log_2\left(\frac{n_i}{N}\right)
\end{equation}

Los posibles valores quedan acotados por:

\begin{equation}
\label{EQ:COTA_ENTROPIA}
0 \leq H(X) \leq 8
\end{equation}

El valor mínimo de la entropía se alcanza cuando todos los bytes del \textit{payload} tienen el mismo valor, es decir, cuando $p(x)=1$. En ese caso, al sustituir en la expresión, cada sumando queda multiplicado por $\log_2(1)=0$, por lo que:

\begin{equation}
\label{EQ:ENTROPIA_MINIMA}
H(X) = 0
\end{equation}

El valor máximo de entropía se alcanza cuando todos los bytes del flujo son equiprobables, es decir, cuando  \begin{equation}p(x_i) = \frac{1}{256}\end{equation}

Sustituyendo en la expresión anterior, se obtiene:

\begin{equation}
\label{EQ:ENTROPIA_MAXIMA}
H(X)
= - \sum_{i=0}^{255} \frac{1}{256}\log_2\left(\frac{1}{256}\right)
= -256 \cdot \frac{1}{256}\log_2\left(\frac{1}{256}\right)
= -\log_2\left(\frac{1}{256}\right)
= 8
\end{equation}

Este cálculo puede realizarse tanto para un paquete individual como para un flujo de paquetes. En este caso, se va a utilizar para ambos, ya que hay algunas características que ahora se explicarán que utilizan la entropía de paquetes en vez de la de flujos. Para obtener la entropía del flujo, simplemente se juntan los bytes de todos los \textit{payload} de los paquetes pertenecientes al mismo flujo y se aplica la expresión anterior.